% A detection reduction, not a computation of the shallow cyclic deformation group.
\begin{lemma}[Proved: cyclic restriction detects the degree-three comparison kernel]
\label{C:S3_cyclic_detection}
Let $A$ be a characteristic-zero DVR of residue characteristic
three containing a primitive cube root $\zeta$, and let $k$ be
a commutative $A$-algebra such that
$K_0(k)\to K_0(k[1/3])$ is injective.
For every
\[
 \xi\in\ker\bigl(R_k(S_3)\longrightarrow R_{k[1/3]}(S_3)\bigr)
\]
one has the exact identity
\[
 \xi=-\operatorname{Ind}_{C_3}^{S_3}
       \bigl([k_\zeta]\operatorname{Res}_{C_3}^{S_3}\xi\bigr).
                                                        \tag{C.1.32}
\]
Consequently restriction to $C_3$ is injective on this kernel.
In particular this applies to
$k=\mathbb Z_3[\zeta]\ltimes\mathbb F_3$.
This is a detection reduction; it does not classify the cyclic
deformation classes.
\end{lemma}
\begin{proof}
Let $K$ be the fraction field of $A$.
The permutation representation induced from $C_2$ has character
equal to the sum of the trivial and standard characters of
$S_3$. The representation induced from the nontrivial character
$\zeta$ of $C_3$ is the standard character. Thus
\[
 1=\operatorname{Ind}_{C_2}^{S_3}(1)
       -\operatorname{Ind}_{C_3}^{S_3}([A_\zeta])
                                                        \tag{C.1.33}
\]
in $R_A(S_3)$, because the generic comparison is injective by
Theorem~\ref{C:DVR}. One can check the character identity at
the identity, a transposition, and a three-cycle: the two
induced characters are respectively $(3,1,0)$ and $(2,0,-1)$.
Every admissible sequence of $A$-lattices splits as a sequence
of $A$-modules. Therefore base change to $k$ preserves the
identity, without requiring that $k$ be flat over $A$.

Since two is a unit, the two idempotents $(1\pm s)/2$ identify
$R_k(C_2)$ with $K_0(k)\oplus K_0(k)$, compatibly with base
change. The assumed injectivity on $K_0$ thus implies
$\operatorname{Res}_{C_2}^{S_3}\xi=0$.
Multiply (C.1.33) by $\xi$ and use the projection formula for
induction. The $C_2$ term vanishes, and the $C_3$ term gives
(C.1.32). This proves the injectivity claim.
For the stated square-zero local ring, both source and target
$K_0$ are $\mathbb Z$ by rank, so the hypothesis holds.
\end{proof}

For clarity, this also eliminates a small apparent obstruction.
In the stated square-zero ring write $\epsilon$ for its residue
ideal generator and set
\[
 s=\begin{pmatrix}1&0\\0&-1\end{pmatrix},\qquad
 t=I_2+\frac{\epsilon}{2}
          \begin{pmatrix}0&1\\1&0\end{pmatrix}.
\]
These matrices satisfy $s^2=t^3=1$ and $sts=t^{-1}$.
The generic representation is the sum of the trivial and sign
lines. Its relative $C_3$ class is zero by
Proposition~\ref{C:cyclotomic_matrix_cancellation}, since its
nilpotent perturbation has trace zero. The lemma therefore shows
that its relative $S_3$ class is also zero. This cancellation
does not cover perturbations of the non-diagonal integral
lattice types in Proposition~\ref{C:cyclotomic_C3_lattices}.
